We now proceed with a detailed presentation of the main research of the proposed project. 


\subsection{Data collection and processing} 
(Goce + Mubbasir)\\
Data cleaning, data collection, sharing, processing, analytics\\
How can we secure the collected data and how to make it available?\\
APIs and libraries for building applications 




\paragraph{Knowledge Representation and Reasoning} 

The processed data from above will be represented as a set of discrete time series, or \emph{fluents} which encode the states and properties of different \avatar which change over time. \todo{give example}. To facilitate efficient reasoning and knowledge discovery in a dynamic smart and inter-connected system of avatars, we will consider a level of abstraction, termed as \emph{events }. Events constitute meaningful interactions between two or more \avatar and serve as the basis for identifying relationships between them towards creating social networks of avatars. \todo{give example}. We define the specific terms and concepts we will use for event-centric knowledge representation and reasoning below. 

\noindent \textbf{Smart Avatars} The notion of smart objects~\cite{Kallmann:1999:DIS:323663.323683} has been popularly used in the graphics and animation community to embed intelligence and semantics in virtual objects. We extend this formalism to represent both IoT objects as well as IoT users. This unified formalism allows us to seamless consider sensors, controllers, actuators, and human users within the same social community.  We define a smart avatar $\smartobj{} \in \world$ as $\smartobj{} = \langle \affordances, \state{} \rangle$ with a set of advertised affordances $\affordances$ and a state \state{}. An affordance $\affordance(\affordanceowner, \affordanceuser) \in \affordances$ is an advertised capability offered by a smart avatar that manipulates the states of the owner of an affordance $\affordanceowner$ and a smart avatar user $\affordanceuser$. 

\noindent \textbf{State.} The state $\state{} = \langle \attributeSymbol, \relationshipSymbol \rangle$ of a smart object $\smartobj{}$ comprises a set of attribute mappings \attributeSymbol (fluents), and a collection of pairwise relationships $\relationshipSymbol$ with other avatars. With this representation, we can make logical inferences between objects using a declarative PROLOG-like knowledge reasoning engine. 
%An attribute \attributemapping{i}{j} is a bit that denotes the value of the $j^{th}$ attribute for $\smartobj{i}$. 



%A specific relationship $\relationship{\cdot}{\cdot}{a}$ is a sparse matrix of $|\world| \times |\world|$, where \relationship{i}{j}{a} is a bit that denotes the current value of the $a^{th}$ relationship between \smartobj{i} and \smartobj{j}. 


%\noindent \textbf{Rules.}  A rule $\Rule_x(i,j) \in \Rules$ between two smart objects \smartobj{i}, \smartobj{j} is true or false, depending on the states and relationships of both objects. Rules allow for logical inference between objects and are used for reasoning such as evaluating whether a character can access a particular room, or manipulate another smart object based on the current world state. They are defined and solved using a declarative PROLOG-like interface 

\noindent \textbf{Events.} Events are used to encode context-specific interactions between two or more smart avatars, and provide an appropriate level of abstraction for knowledge discovery  An event is formally defined as $\event{} = \langle \tree,  \precondition{}, \postcondition{} \rangle$. A precondition $\precondition{}: \groupState{\smartObjectGroup{}} \leftarrow \{\true, \false \} $ is a logical expression on the compound state \groupState{\smartObjectGroup{}} of a particular set of smart avatars $\smartObjectGroup{}: \{ \smartobj{1}, \smartobj{2}, \ldots \smartobj{|\roleSet{}|} \}$ that checks the validity of the states of each smart object. \precondition{} is represented as a conjunction of clauses $\clause{}{} \in \precondition{}$ where each clause \clause{}{} is a literal that specifies the desired attributes of smart objects, and relationships between pairs of participants. A precondition is fulfilled by $\smartObjectGroup{} \subseteq \world $ if $\preconditionFunction{\event}{\smartObjectGroup{}} = \true$. The event postcondition $\postcondition{} : \groupState{} \rightarrow \nextGroupState{}$ transforms the current state of all event participants \groupState{} to \nextGroupState{} by executing the effects of the event. We can extend this definition to model non-deterministic, fuzzy events with a probabilistic notion of success or failure. Events may optionally have a controller which defines the series of affordance activations within the smart avatars to produce its desired outcome. We represent this control logic using an extended version of Behavior Trees that facilitate parameterization. Parameterized Behavior Tree (PBT)~\cite{Shoulson:2011:PBT:2177817.2177835} are an effective model for representing coordinated control logic  between multiple smart avatars. 

PI Kapadia has extensive prior experience in developing event-centric knowledge bases for inference and reasoning in virtual worlds~\cite{Shoulson:2013:EPA:2522628.2522629,2015-fdg-bstl,Kapadia:2015:CAI:2699276.2699279,Kapadia:2016:CCN:2982818.2982846} which will be extended to represent the IoT domain. 


%ADAPT~\cite{Shoulson:2013:EPA:2522628.2522629}
%BSTL~\cite{2015-fdg-bstl,Kapadia:2015:CAI:2699276.2699279}
%behavior authoring~\cite{behaviorCGA,Kapadia:2011:BAC:1944745.1944779}
%story world~\cite{AIIDE1511583}
%CANVAS~\cite{Kapadia:2016:CCN:2982818.2982846}





%$\roleSet{} = \{ \role{i} \}$  define the desired roles for each participant. \role{i} is a logical formula specifying the desired value of the immutable attributes \attributemapping{\cdot}{j} for \smartobj{j} to be considered as a valid candidate for that particular role in the event.  


%An event instance $\instance = \langle \event, \smartObjectGroup{} \rangle$ is an event \event populated with an ordered list of smart object participants \smartObjectGroup{}.  $\preconditionFunction{\event}{\smartObjectGroup{}} = \true$. The event postcondition $\postcondition{} : \groupState{} \rightarrow \nextGroupState{}$ transforms the current state of all event participants \groupState{} to \nextGroupState{} by executing the effects of the event. When an event fails, $\nextGroupState{} = \groupState{}$. An event instance $\instance = \langle \event, \smartObjectGroup{} \rangle$ is an event \event populated with an ordered list of smart object participants \smartObjectGroup{}. 

%? which represent the specific values of avatar fluents at specific points. Events occur on a set of avatars which are the participants, and produce a change in state of these avatars which are defined as event postconditions. Events may additionally be pre-conditioned on the participant states to indicate what conditions must be satisfied in order for an event to successfully execute. 

This symbolic representation of avatars and theiry dynamic states in terms of key events will facilitate the development of an event calculus ? allowing us make inferences about the relationships and properties of these avatars using first and second-order logic. This reasoning will be used as the basis for knowledge discovery described in Section XX. Knowledge discovery will entail the identification of relationships between the properties of different avatars, clustering of avatar to create compound entities, and the discovery of salient events. 


 

\subsection{Knowledge discovery: social network of avatars (virtual world)}
(Ashfaq + Vladimir)

Tasks:
\begin{itemize}
	\item Build spatio-temporal avatar profiles
	\item Group avatars into groups according to similar profiles. I.e., construct social network of profiles and do graph clustering.
	\item Predict avatar behavior from group profiles
\end{itemize}


We will use Gaussian Process dynamic models (GPDM) as profiles. GPDMs can be used for both modeling isolated avatars and interactions between avatars.  We can use a framework similar to \cite{Shen2015-jr,Shen2015-ft,Shen2012-vi}.  Then extend using our trajectory refinement and optimization approaches \cite{yoon2016}.

However, we need to consider the fact that this system is heterogeneous, i.e., things and humans.  So should have a hierarchical model, one level for things, one for humans, then merge on higher level.


\subsection{Security, privacy and anonymity}



Security poses a great challenge for IoT. In addition to security threats common for networked systems in general, IoT systems are further confronted by the fact that they necessarily rely on wireless signals for communication. We propose to address challenges stemming from the \emph{physical layer} attacks on IoT systems. In particular, while Multi-User Multiple Input Multiple Output (MU-MIMO) systems are considered the ultimate next step in the wireless bandwidth race (because they leverage spatial multiplexing to send independent data streams to multiple terminals simultaneously, effectively improving spatial reuse), they could be utilized as an effective and sophisticated \emph{denial-of-service} tool. Indeed, by creating controlled signals in an area, it can selectively disrupt performance of Wi-Fi networks and sensors in an area, while still making it hard to detect or classify such behavior (precisely due to uneven signal properties at different locations). Such an attack is relevant in many scenarios. Consider a location densely populated with wireless APs. Can one of the APs turn itself into a “rogue” mode to disrupt the other APs, i.e., “clear the air”, and thus improve performance for itself, while staying largely “invisible” to the other APs? Another example is Wi-Fi backscatter, \emph{i.e.}, scenarios where a device backscatters Wi-Fi signals to communicate. In presence of the attack, such communication might be severely degraded, while at the same time it becomes hard to know that the attack is taking place. Given that Wi-Fi backscatter was recently proposed as a solution to communication with implanted medical devices, this scenario significantly raises the bar for addressing this emerging problem. Below, we first outline our current results in this domain, and the outline a research agenda. 


{\bf Background and Current Work}

A fundamental requirement in MU-MIMO systems is that up-to-date Channel State Information (CSI), as perceived by the clients, must be obtained by the AP over short time scales. We show that both explicit and implicit \emph{client-fed} CSI feedback open the door to a fundamental vulnerability of MU-MIMO systems. We demonstrate that a single attacker (a Neo) can mount strategic attack, called Matrix-bending, that can dramatically reduce the throughput
of an MU-MIMO system. Moreover, we show that a small number of Neos (or a single Sybil-enabled Neo) can reduce the throughput of an MU-MIMO system to a level \emph{below} that of a corresponding single-user (SU-MIMO) system. Such a degradation translates to a significantly impaired user experience, \emph{e.g.}, increasing the number of perceivable freeze periods while video-streaming by a factor of $10$ and their total duration by a factor of $20$,
spending $34\%$ of the streaming duration frozen, in comparison to merely $1.7\%$ when using MU-MIMO under the same conditions~\cite{mu-mimoUXeffect,mimo13sigcomm}.

The key behind Matrix-bending is a strategically disruptive behavior in which a Neo reports its CSI with the goal of deliberately ``populating'' the clients' signal space as seen by the AP. By doing so, the Neo confuses the AP, which in turn fails to create ``good'' user clusters. In particular, by becoming a desired member of most client groups, and by driving the system towards generating extremely sub-optimal groups, the Neo effectively breaks client selection algorithms, leading to substantial throughput degradation


We demonstrate the following important features of the Matrix-bending attack:
($i$) Matrix-bending is \emph{a low-profile} attack. A Neo does not apply any jamming, and transmits data using only a single antenna. Moreover, it generates CSI that is statistically indistinguishable from the CSI generated by regular clients, and is capable of achieving its full impact while generating traffic at volumes that are more than an order of magnitude smaller than those generated by a regular client.
($ii$) A single Neo has the capacity to considerably decrease MU-MIMO's throughput gains, \emph{i.e}, by up to one third, even in the presence of $100$ legitimate clients. To bring an MU-MIMO system performance to the level of a corresponding single-user system, the number of Neos in the system need not be large, and does not depend on the number of regular clients in the system. Instead, the required number of Neos corresponds to the number of clients in a beamforming group, which is theoretically upper bounded by the number of antennas at an AP, and in practice often bounded by a smaller number, \emph{e.g.}, $4$ in 802.11ac~\cite{std11ac}.

We design and implement a real-time 802.11ac beamforming system on a software-defined radio 
platform that allows us to evaluate Matrix-bending on over-the-air 5~GHz channels with complete access to CSI~\cite{warpProject}. We evaluate MU-MIMO client selection 
algorithms and policies in the presence of Neos. We show that falsifying the CSI, \emph{i.e.}, the V-matrix, of a single Neo is sufficient to reduce a substantial portion of MU-MIMO's throughput gain, \emph{i.e.}, approximately by one third. Using $3$ Neos, the throughput degrades by approximately $3\times$, often experiencing performance below
the corresponding SU-MIMO system. We find that Matrix-bending's effect is most dominant over APs that attempt to maximize throughput,  while it is expectedly less pronounced for the schemes that do not focus on throughput.

We further study the potential defenses against Matrix-bending.  First, we evaluate random user selection algorithms. While they are less vulnerable to Matrix-bending attacks than other algorithms, such a feature comes with a significant cost: the throughput of such systems decreases below the levels achieved in SU-MIMO systems, even in the absence of a Matrix-bending attack. Second, we evaluate client CSI encryption. We find that Matrix-bending is quite effective in such scenarios when it has \emph{no knowledge about the CSI from other clients.} The key reason is that as long as the Neos' responses are near-orthogonal to each other, they will affect other users as well. Third, we find Matrix-bending to be equally effective when it has no knowledge about its own CSI. In particular, we evaluate a method that utilizes unknown sounding preambles which aim to prevent CSI manipulation. However, we show that all realistic preamble transformations still preserve the \emph{relative signal properties} such that the attack is still effective.

{\bf Research Agenda}

We propose to address challenges stemming from the \emph{physical layer} attacks on IoT systems. We will evaluate the attacks, and defenses, enabled by MU-MIMO systems. In particular, we propose research in the following directions:

\emph{Silent DoS.} In this scenario, the attacker can jam all communication between clients and an AP in both directions, while \emph{quiet-bubbling} the AP’s antennas, therefore preventing it from sensing the disturbance. The key research questions are how to effectively use \emph{implicit} CSI in order to quiet bubble the APs and to understand the impact on the AP’s performance.

\emph{Silent scrambler.} In this scenario, we consider Wi-Fi-enabled backscatter communication, where signals generated by a Wi-Fi source are backscattered by a client device. In this case, the attacker node (an 8-antenna 802.11ac AP) transmits a time-varying signal while quiet-bubbling the Wi-Fi source. The backscatter client-device responds to the scrambled signal it receives, thus disrupting the communication.
\emph{Power-killer attack.} In this scenario, we again assume the above Wi-Fi-enabled backscatter communication. The goal of this attack would be to null the signal at the backscatter device, such that the energy harvested by the backscatter client is insufficient to enable transmissions. The key research question is to determine critical time-scales and its impact on the effectiveness of the attack.

\emph{Countermeasures.} We propose to analyze countermeasures to the above attacks. The most promising approach for the silent DoS attack is to equip APs with an additional, passive antenna, that would collect the information from the environment. Then, we would further conduct statistical analysis to find critical time scales that enable the attack detection with appropriate confidence guarantees. For the silent-scrambler and power-killer attacks, we will pursue solutions that do not require additional hardware. In all cases, we will conduct experiments on a real-time 802.11ac beamforming system on a software-defined radio 
platform on over-the-air 5~GHz channels~\cite{warpProject}.



(Farid + Alex)\\
\\
Data security\\
System level security\\
Privacy, anonymity\\
\\
Greedy behavior (?)\\
Key sharing problem\\ 
Privacy, trust and anonymity \\

\subsection{Implementation and validation }
(All)\\
\\
Letters of Collaboration)\\
Campus wide implementation \\
Letters of Commitment\\
Shall we rely on KAA?\\
APIs: flexibility of formats/schemas\\

Given the interdisciplinary nature of the tasks in the proposed project, it is necessary to have well-defined mechanisms for assessing the validity and benefits of our findings. We will base our demonstration on variations of the two methodologically interdependent scenarios used in the introductory section.  We now discuss both the broad categories of evaluation environments, and the details of the particular evaluation methodologies. 
We will adopt a three-fold analysis of the different aspects of our system: (1) agent-based computer simulations, (2) experiments with real humans controlling avatars in shared virtual worlds, and (3) experiments with real humans in actual real-life environments. Each of these experimental paradigms allow us stress our system along a combination of the following three axes: (1) human stress, (2) replicating the noise inherent in using humans as sensors, and (3) scalability ? i.e., the number of participants (humans or agents). Agent-based simulations harness computation to simulate high-stress events in arbitrarily large crowds of computer agents, but are unable to accurately model humans. Shared online virtual worlds provide a unique opportunity to conduct experiments with real humans in controlled settings. While experiments with human subjects in the real world offer the closest possible parallel to reality, these experiments are limited in number of subjects and the extent of stress that we can induce on the participants. However, by leveraging combinations of all 3 paradigms, we are uniquely positioned to evaluate our models and technologies in a rigorous manner. 


\begin{wrapfigure}{R}{0.30\textwidth} \vspace{-3mm}
	\centerline{\includegraphics[width=0.30\textwidth]{evaluation-plan.png}}
	\vspace{-3mm} \caption{\small Evaluation using Agent simulations, Shared Online Virtual Worlds, and Real World Experiments.}
	\label{fig1}
	\vspace{-3mm}
\end{wrapfigure}

\subsubsection{Simulation-based Evaluations: Agents. }

The first part of our efforts will focus on developing simulator for the two demonstration settings. Throughout the different phases of the project we will be capitalizing upon ? as well as augmenting the current capabilities of ? the SIDnet SWANS simulator [195] and SteerSuite: The Human Movement and Crowd Simulator [196], being developed at Northwestern and Rutgers respectively. SIDnet is used to evaluate various routing and tracking protocols, shapes detection, and even to evaluate the benefits of high-level programming constructs for WSN-users [113, 160, 187, 191]. SteerSuite is an open-source framework for simulating human movement, deliberation, and is especially geared towards dense crowd situations. We will integrate the functionality of these two systems and further augment them with the capability of having heterogeneous nodes and capturing various mobility models and communications between subsets of such nodes. In addition, we will augment its visualization component as well as introduce a new class of nodes with the proper interfaces to encode the following important features: (1) camera/video/audio based sensing and (2) human and human-operated sensors.  For the multimodal mashup scenario, we will simulate a simplified version of the ?Gold Miner? using computer imagery, with virtual objects and sensors. 
The main use of the simulator in the initial stages of the proposal will be to test the various impacts of uncertainty (sensors, network structure/communication) on the quality of detecting the novel categories of spatio-temporal predicates (cf. Section 3.4) and the impact of different policies for local state-estimation and information extraction (cf. Section 3.1 and 3.3). We will employ a large variety of crowd datasets to train and validate our computational models, and in particular leverage our collaboration with Prof. Dirk Helbing who has significant expertise in crowd data collection, analysis, and modeling [197].  In particular, we will integrate our simulation infrastructure with the NervousNet platform [198] that is being actively developed by Prof. Helbing and his colleagues as part of the larger effort on the Planetary Nervous System project [199], which provides a large-scale distributed research platform for real-time mining of social activities using heterogeneous sensor networks.


\subsubsection{Task E-2. Game Playing Evaluation: Avatars. }

Agent-based simulations, such as those described above provide an efficient way to test and evaluate the different aspects of our system in an automated fashion. However, in an effort to come closer to actual human experimentation, we will conduct experiments using real human users in shared online virtual environments. The Rutgers team is uniquely positioned to conduct these experiments ? Co-PI Kapadia has developed HeapCraft [200, 2001]: a modular, extensible, and open framework for studying human behavior in shared virtual worlds such as Minecraft. HeapCraft has already been successfully demonstrated to study player behavior and incentivize cooperation in online societies [202, 203]. Additionally, Kapadia and his collaborators have also conducted lab experiments to study crowd evacuations using serious games.
We will augment the existing HeapCraft framework to design the two scenarios described in Task E-3 below. HeapCraft?s modular and extensible nature allows us to easily develop emulators for the technologies being developed here, which can be rigorously and robustly tested in shared virtual worlds using real human users. In addition, we will augment the knowledge-base for emulating behavioral aspects of the avatars using large datasets of preferences and semantic relations from real social networks obtained from 4C Insights Inc. (Letter of Commitment is available in the Supplementary Documents). 
Our experiments will be conducted in two phases: (I) laboratory studies where up to 40 participants will be recruited to participate in a controlled experiment; (II) online server with our tools released to setup a virtual world with hundreds of human-users  expected to participate. 


\subsubsection{Task E-3. Real-World Environment Evaluations.  }

XX

 
